\section{本讲小结}

光路计算从交点和当地法线开始。反射保留传播矢量的切向分量并反转法向分量，折射保留 $n\sin\theta$。球面、棱镜不同面的法线通常不同，上一段的入射角不能原样带到下一面。全反射还要求由较大折射率射向较小折射率，并严格超过临界角；等号对应掠射边界。

镜像展开保持长度和交角，用辅助直线求反射路径后，应折回并逐次检查真实碰镜和穿屏位置。分层介质中，各层以同一参数 $p$ 连接；横向投影、法向厚度、实际路程分别进入端点约束、几何关系和计时。连续分层的转向点由实方向边界决定，不能当作均匀介质内的一条折线。

费马原理要求允许变化下的光程驻值；最小性、交点范围和传播支路应由具体模型另行检验。等光程设计须从同一入射波前出发，给各介质路程分别加权，并用局部反射或折射定律核对法线方向。近轴方向的一阶近似与光程差所需的二阶长度项，应按最终所求量分别保留。
